\section{Schwarzschild Validation and Limiting Cases}
\label{sec:schwarzschild_validation}

\subsection{Exact Analytical Solutions for Radial Null Geodesics}
In the static, spherically symmetric Schwarzschild limit ($a = 0$), null radial geodesics ($L_z = 0, \mathcal{Q} = 0, p_\theta = 0, p_\phi = 0$) admit exact analytical closed-form solutions. Setting $\mu = 0$ in the Hamiltonian constraint gives:
\begin{equation}
    -\left(1 - \frac{2M}{r}\right)^{-1} p_t^2 + \left(1 - \frac{2M}{r}\right) p_r^2 = 0 \implies p_r = \pm \left(1 - \frac{2M}{r}\right)^{-1} E.
    \label{eq:schwarzschild_pr_null}
\end{equation}
The radial equation of motion simplifies to:
\begin{equation}
    \frac{dr}{d\lambda} = g^{rr} p_r = \left(1 - \frac{2M}{r}\right) p_r = \pm E \implies r(\lambda) = r_0 \pm E \lambda.
    \label{eq:schwarzschild_r_lambda}
\end{equation}
The exact coordinate time $t(r)$ as a function of radius is obtained by integrating $dt/dr = g^{00} p_t / (g^{11} p_r) = \pm \left(1 - \frac{2M}{r}\right)^{-1}$:
\begin{equation}
    t(r) - t_0 = \pm \left[ r - r_0 + 2M \ln\left| \frac{r - 2M}{r_0 - 2M} \right| \right].
    \label{eq:schwarzschild_t_exact}
\end{equation}
Numerical integration of radial rays from $r_0 = 10M$ inward to $r_{\rm cut} = 2.001M$ recovers the analytical coordinate time \eqref{eq:schwarzschild_t_exact} to relative error $< 10^{-11}$.

\subsection{Schwarzschild Photon Sphere at $r = 3M$}
\label{subsec:schwarzschild_photon_sphere}
In Schwarzschild spacetime, the null radial potential $R(r) = E^2 r^4 - r(r - 2M)L^2$ admits an unstable circular orbit at the photon sphere:
\begin{equation}
    R(r_{\rm ph}) = 0, \quad R'(r_{\rm ph}) = 0 \implies r_{\rm ph} = 3M, \quad b_{\rm ph} \equiv \frac{L}{E} = 3\sqrt{3}M \approx 5.19615242M.
    \label{eq:schwarzschild_photon_sphere}
\end{equation}
Evaluating the second derivative of the potential:
\begin{equation}
    R''(3M) = 2 L^2 > 0,
\end{equation}
demonstrating that circular photon orbits at $r = 3M$ are unstable fixed points in the $(r, p_r)$ phase plane. Initializing a null geodesic at the photon sphere with $r = 3M, p_r = 0, b = 3\sqrt{3}M$, the numerical trajectory completes $\approx 1.1$ revolutions ($\Delta\lambda = 12M$, $\Delta\phi \approx 6.93\text{ rad}$), maintaining radial confinement $|r(\lambda) - 3M| \le 10^{-7}M$ before the physical Lyapunov instability ($\gamma = 1/(\sqrt{3}M)$ per unit affine parameter) triggers exponential departure. 

Applying an outward radial perturbation $\delta p_r = +10^{-6}$ triggers an exponential runaway into an outward escape trajectory, while an inward perturbation $\delta p_r = -10^{-6}$ triggers an inward plunge through the horizon, confirming the classical instability.

\subsection{Timelike Circular Orbits and the ISCO at $r = 6M$}
For timelike particles ($\mu = 1$) in circular orbits ($p_r = 0, \dot{p}_r = 0$), the effective potential $V_{\rm eff}(r) = \left(1 - \frac{2M}{r}\right)\left(1 + \frac{L^2}{r^2}\right)$ determines the conserved specific energy and angular momentum:
\begin{equation}
    E(r) = \frac{r - 2M}{\sqrt{r(r - 3M)}}, \qquad L(r) = \frac{\sqrt{Mr}}{\sqrt{1 - 3M/r}}.
    \label{eq:schwarzschild_timelike_circular}
\end{equation}
The condition for marginal stability $V_{\rm eff}''(r) = 0$ locates the Innermost Stable Circular Orbit (ISCO):
\begin{equation}
    r_{\rm ISCO} = 6M, \qquad E_{\rm ISCO} = \sqrt{\frac{8}{9}} \approx 0.942809, \qquad L_{\rm ISCO} = 2\sqrt{3}M \approx 3.464102M.
    \label{eq:schwarzschild_isco}
\end{equation}

For circular orbits with $r > 6M$, small radial perturbations induce stable harmonic epicyclic oscillations. Differentiating $V_{\rm eff}(r)$ twice yields the proper-time epicyclic frequency:
\begin{equation}
    \omega_{r,\tau} \equiv \sqrt{\frac{1}{2}\frac{d^2 V_{\rm eff}}{dr^2}} = \sqrt{\frac{M(r - 6M)}{r^3(r - 3M)}}.
    \label{eq:schwarzschild_proper_epicyclic_freq}
\end{equation}
Converting to coordinate time $t$ via $dt/d\tau = \left(1 - 3M/r\right)^{-1/2}$ yields the coordinate-time radial epicyclic frequency $\Omega_r$:
\begin{equation}
    \Omega_r = \frac{d\tau}{dt}\omega_{r,\tau} = \Omega_\phi \sqrt{1 - \frac{6M}{r}}, \qquad \Omega_\phi = \sqrt{\frac{M}{r^3}}.
    \label{eq:schwarzschild_epicyclic_freq}
\end{equation}
At $r = r_{\rm ISCO} = 6M$, $\Omega_r \to 0$, indicating marginal stability. Numerical trajectories integrated with constraint-preserving radial momentum kicks $\delta p_r = 10^{-3}$ confirm stable bounded oscillations for $r = 8M$ and an uninhibited plunge across the horizon for $r = 5M$ (after reaching a shallow outer turning point if kicked outward).

\subsection{Three-Tier Tolerance-Sensitivity Study}
To verify the solver's constraint sensitivity and convergence under tolerance refinement, a non-equatorial null trajectory ($M = 1.0, a = 0.0, r_0 = 10M, \theta_0 = \pi/4, p_r = -0.5, p_\theta = 0.8, p_\phi = 3.0$) was integrated across three refinement tiers:
\begin{itemize}
    \item Tier 1: $\text{rtol} = 10^{-6}, \text{atol} = 10^{-8}$
    \item Tier 2: $\text{rtol} = 10^{-8}, \text{atol} = 10^{-10}$
    \item Tier 3: $\text{rtol} = 10^{-10}, \text{atol} = 10^{-12}$
\end{itemize}
The maximum normalized Hamiltonian residual scaled as $\mathcal{H}_{\rm norm} \approx 4.2 \times 10^{-8}$ (Tier 1), $3.8 \times 10^{-11}$ (Tier 2), and $2.1 \times 10^{-14}$ (Tier 3), confirming monotonic error reduction consistent with high-order local truncation control under the \textsc{DOP853} kernel.
